% !PRINTED 623
\weberpage{623}
% No page number is visible within the captured band crop (the top
% band is blank above the Abschnitt title, as this chapter opener carries no
% running head, and the foot of the page below the footnote is blank too). 623 is inferred from direct pagination continuity: this page's footnote
% 1) is continued verbatim at the head of the footnote apparatus on the next page,
% whose own running head clearly prints "624".

\begin{center}
\gesperrt{Twenty-fourth Chapter.}

\medskip
{\Large Algebraic Functions of One Variable.}
\end{center}

\rule{5em}{0.4pt}

\begin{center}
§ 170. \gesperrt{Introductory.}
\end{center}

The theory of algebraic functions of one variable is the starting point
of the general investigations of \gesperrt{Abel} on the new
transcendents, which have since received the name \glqq \gesperrt{Abelian functions}\grqq{} and which are the highest generalisation of the
elliptic functions. The principal problems of this theory have been
brought to a certain conclusion by the papers of \gesperrt{Riemann},
\gesperrt{Weierstrass}, \gesperrt{Clebsch}, \gesperrt{Brill} and
\gesperrt{Noether}. In particular, in the conception of many-sheeted
(\gesperrt{Riemann}) surfaces \gesperrt{Riemann} has created a means for the
investigation of these functions that is extraordinarily effective
through its intuitive clarity.

All these investigations, however, which rest on the one hand upon the
theory of functions, on the other hand, as with \gesperrt{Clebsch},
\gesperrt{Brill}, \gesperrt{Noether}, upon the methods of rational
algebra (theory of forms and invariants, or of algebraic curves), must
always make certain restrictions; they must exclude certain
\glqq exceptional cases\grqq{} and content themselves with the treatment
of the so-called general cases. Not subject to this restriction is the
theory of algebraic functions worked out on the analogy of number theory
by \gesperrt{Dedekind} and myself, of which a survey is to be given
here\textsuperscript{1)}.

\smallskip
\noindent\rule{5em}{0.4pt}

\noindent\textsuperscript{1)}~On the literature concerning algebraic and
\gesperrt{Abel}ian functions let the following be mentioned here: \gesperrt{Abel},
Mém. sur une propriété générale d'une classe très-étendue de fonctions
transcendentes (presented to the Paris Academy in 1826, Werke, new
% [sic] printed "transcendentes", Abel's French title (and French spelling) has "transcendantes"
edition by \gesperrt{Sylow} and \gesperrt{Lie}, Bd.\,I, S.\,145).
\gesperrt{Riemann}, Theorie der \gesperrt{Abel}schen Funktionen, Crelles Journ.,
Bd.\,54, 1857, Werke 2. Aufl., Nr.\,88. \gesperrt{Weierstrass},
Vorlesungen 1875/76, Werke, Bd.\,IV. \gesperrt{Clebsch}, Über die
%CONT>>
